\documentclass[10pt]{amsart}
\usepackage[a4paper,margin=23mm]{geometry}
\usepackage{amsmath,amssymb,amsthm,mathtools,hyperref,graphicx,xspace,letltxmacro,xpatch}
\usepackage[english]{babel}
\usepackage{xfrac}

\usepackage{mathrsfs}
\let\mathcal\mathscr

\newcommand{\ldbrack}{[\![}
\newcommand{\rdbrack}{]\!]}
\def\vcentcolon{\mathrel{\mathop\ordinarycolon}}
\providecommand\ordinarycolon{:}
\providecommand*\coloneqq{\vcentcolon\mathrel{\mkern-1.2mu}=}
\providecommand*\eqqcolon{=\mathrel{\mkern-1.2mu}\vcentcolon}
\let\eqcolon\eqqcolon
\newcommand\mto{\mathchoice{\longmapsto}{\mapsto}{\mapsto}{\mapsto}}
\newcommand\implique{\mathchoice{\Longrightarrow}{\Rightarrow}{\Rightarrow}{\Rightarrow}}
\newcommand{\RedefinitSymbole}[1]{%
\expandafter\let\csname old\string#1\endcsname=#1
\let#1=\relax
\newcommand{#1}{\csname old\string#1\endcsname\,}%
}
\RedefinitSymbole{\forall} \RedefinitSymbole{\exists}

\usepackage{nicefrac}

\newcommand{\underdot}[1]{\smash{\underset{\raisebox{3pt}{.}}{#1}}}

\usepackage{tikz}
\usetikzlibrary{arrows}
\usepackage{tikz-cd}
\tikzcdset{every label/.append style = {font = \normalsize}}
\usepackage[capitalise]{cleveref}
\crefdefaultlabelformat{#2\textup{#1}#3}%

\theoremstyle{plain}
\newtheorem{theorem}{Theorem}[section]
\newtheorem{corollary}[theorem]{Corollary}
\newtheorem{proposition}[theorem]{Proposition}
\newtheorem{lemma}[theorem]{Lemma}

\theoremstyle{definition}
\newtheorem{example}[theorem]{Example}
\newtheorem{remark}[theorem]{Remark}
\newtheorem{question}[theorem]{Question}
\newtheorem{claim}[theorem]{Claim}

\Crefname{section}{Section}{Sections}
\Crefname{subsection}{Section}{Sections}
\Crefname{figure}{Figure}{Figures}

\AtBeginDocument{%
\renewcommand{\theenumi}{\alph{enumi}}%
\renewcommand{\theenumii}{\arabic{enumi}}%
}

\newcommand{\tribar}[1]{\mathopen{| {\kern -1.5pt} | {\kern -1.5pt} |} {#1}
\mathclose{| {\kern -1.5pt} | {\kern -1.5pt} |}}

\newcommand{\RR}{\mathbb{R}}
\newcommand{\ZZ}{\mathbb{Z}}
\newcommand{\NN}{\mathbb{N}}
\newcommand{\EE}{\mathbb{E}}
\newcommand{\FF}{\mathbb{F}}
\newcommand{\GG}{\mathbb{G}}
\newcommand{\KK}{\mathbb{K}}
\newcommand{\MM}{\mathbb{M}}
\newcommand{\PP}{\mathbb{P}}
\newcommand{\WW}{\mathbb{W}}
\newcommand{\id}{\mathrm{id}}
\newcommand{\cB}{\mathcal{B}}
\newcommand{\cE}{\mathcal{E}}
\newcommand{\cG}{\mathcal{G}}
\newcommand{\cI}{\mathcal{I}}
\newcommand{\cM}{\mathcal{M}}
\newcommand{\cR}{\mathcal{R}}
\newcommand{\cS}{\mathcal{S}}
\newcommand{\cU}{\mathcal{U}}
\newcommand{\st}{\;\mathord{;}\;}
\DeclareMathOperator{\supp}{supp}

\newcommand*\circled[1]{\tikz[baseline=(char.base)]{
\node[shape=circle,draw,inner sep=1pt] (char) {\footnotesize{#1}};}}

\newcommand{\GL}{\mathit{GL}}
\newcommand{\SL}{\mathit{SL}}
\DeclareMathOperator{\vol}{vol}
\newcommand{\WEDGE}{\mathsf{\Lambda}}
\newcommand{\dd}{\mathrm{d}}
\newcommand{\uu}{\mathrm{u}}
\renewcommand{\ss}{\mathrm{s}}
\newcommand{\loc}{\mathrm{loc}}
\newcommand{\Aut}{\mathrm{Aut}}
\newcommand{\End}{\mathrm{End}}
\DeclareMathOperator{\bol}{bol}
\newcommand{\Gr}{\mathcal{G}}
\DeclareMathOperator*{\linf}{linf}

\let\epsilon\varepsilon
\let\phi\varphi
\let\angle\measuredangle


\usepackage{xfrac,xcolor,graphicx}\graphicspath{{./figures/}}\providecommand{\skpt}{\par\smallskip\noindent}\providecommand{\pointir}{.\ }\providecommand{\eg}{e.g.}\providecommand{\ie}{i.e.}\providecommand{\resp}{resp.}\providecommand{\cf}{cf.}\providecommand{\etc}{etc.}\let\Subsection\subsection\let\Subsubsection\subsubsection
